% Conclusion: Structural stagnation detection in LLM coding agents (EMNLP 2026 ARR)
\section{Conclusion}
\label{sec:conclusion}

We presented a systematic empirical study of progress stagnation detection in LLM coding agents, comparing structural features, neural embedding probes, and LLM-as-judge approaches against silver labels generated via weak supervision.
Three findings emerge.

First, stagnation signals are fundamentally structural.
Thirty hand-crafted features achieve 0.9951 in-domain and 0.967 cross-agent macro-F1, decisively outperforming a GRU probe (+8.96\,pp in-domain) and an LLM judge (+33.3\,pp)---all without GPU and in under 5\,ms per window.
This result holds across three embedding encoders (Qwen3, BGE-M3, E5-large), each yielding probe F1 above 0.90, confirming that the structural advantage is independent of encoder choice.
The practical implication is immediate: real-time agent monitoring need not rely on expensive representation learning or LLM inference when lightweight structural features transfer across agent architectures with minimal degradation.

Second, labeling functions exhibit systematic but agent-specific temporal dynamics.
Command repetition leads stagnation onset by 1.87 steps on SWE-Agent but reverses to a lagging role on OpenHands (+3.32 steps), where its 92.3\% firing rate and 71\% false positive rate corrupt silver labels---removing it improves fine-tuned probe F1 by 17.6\,pp.
This finding has methodological significance for weak supervision on sequential data: labeling functions validated on one domain can become adversarial on another, necessitating per-agent validation and the circular evaluation methodology we develop to surface such pathologies.

Third, temporal sophistication in label aggregation does not improve detection.
Three temporal label models---TWV, HMM, and lag-aware HMM---all fail to surpass majority vote, with TWV actually degrading downstream performance by 2.57\,pp.
This negative result constrains future research: in the low-conflict regime of 3 binary LFs (95.1\% agreement at conf${\geq}$2), temporal alignment introduces more noise than signal.

These results suggest several directions for future work.
Expanding the labeling function set beyond three functions could create the inter-LF conflict necessary for temporal aggregation to add value.
Extending evaluation to non-coding agent tasks (web browsing, data analysis, multi-agent collaboration) would test whether structural dominance generalizes beyond software engineering.
Most critically, evaluation against human-annotated ground truth---for which we have prepared a 650-window gold set with a validated annotation pipeline---will determine whether structural features detect \emph{real} stagnation or merely reconstruct labeling function outputs.
